\documentclass[aspectratio=169]{beamer}

% Tema y configuración
\usetheme{Madrid}
\usecolortheme{default}

% Paquetes necesarios
\usepackage[utf8]{inputenc}
\usepackage[spanish]{babel}
\usepackage{amsmath}
\usepackage{amsfonts}
\usepackage{amssymb}
\usepackage{graphicx}
\usepackage{listings}
% Habilitar UTF-8 dentro de lstlisting (acentos, ñ, etc.)
\usepackage{listingsutf8}
\usepackage{xcolor}
\usepackage{verbatim}
% Para dibujos con TikZ
\usepackage{tikz}
\usetikzlibrary{positioning,trees}

% Configuración de listings para código
\lstset{
    basicstyle=\ttfamily\footnotesize,
    keywordstyle=\color{blue}\bfseries,
    commentstyle=\color{green},
    stringstyle=\color{red},
    showstringspaces=false,
    breaklines=true,
    frame=single,
    numbers=left,
    numberstyle=\tiny\color{gray},
    inputencoding=utf8,
    % Mapeo de caracteres españoles para que no falle la compilación
    literate={á}{{\'a}}1 {é}{{\'e}}1 {í}{{\'\i}}1 {ó}{{\'o}}1 {ú}{{\'u}}1
             {Á}{{\'A}}1 {É}{{\'E}}1 {Í}{{\'I}}1 {Ó}{{\'O}}1 {Ú}{{\'U}}1
             {ñ}{{\~n}}1 {Ñ}{{\~N}}1 {ü}{{\"u}}1 {Ü}{{\"U}}1
}

% Información del documento
\title{Hoja 4: Ejercicio 24}
\subtitle{Ramificación y poda}
\author{Diego Rodríguez Cubero}
\institute{UCM}

% Logo (opcional)
% \logo{\includegraphics[height=1cm]{images/logo.png}}

\begin{document}

% Diapositiva de título
\begin{frame}
    \titlepage
\end{frame}

% Tabla de contenidos
\begin{frame}{Contenidos}
    \tableofcontents
\end{frame}

% Sección 1
\section{Enunciado del problema}
\begin{frame}{Enunciado del problema}
    \begin{block}{Ejercicio 24}
        El problema \textbf{ISO\_SUBGRAPH} consiste, dados dos grafos
        $G = (V_G, A_G)$ y $H = (V_H, A_H)$, en determinar si $G$ tiene un
        subgrafo isomorfo a $H$.

        Se pide demostrar que \textbf{ISO\_SUBGRAPH} es NP-completo, sabiendo
        que el problema de decision \textbf{CLIQUE} lo es.

        Decir que $G$ tiene un subgrafo isomorfo a $H$ significa que existe un
        grafo $G' = (V_{G'}, A_{G'})$ subgrafo de $G$, es decir:

        \[
        V_{G'} \subseteq V_G, \qquad
        A_{G'} = A_G \cap (V_{G'} \times V_{G'})
        \]

        y una biyeccion $f : V_{G'} \to V_H$ tal que:

        \[
        (u, v) \in A_{G'} \Longleftrightarrow (f(u), f(v)) \in A_H
        \]
    \end{block}
\end{frame}

% Sección 2
\section{Demostración de NP-completitud}
\begin{frame}{Demostración de NP-completitud}
    Para demostrar que \textbf{ISO\_SUBGRAPH} es NP-completo, debemos
    mostrar dos cosas:

    \begin{itemize}
        \item \textbf{ISO\_SUBGRAPH} está en NP.
        \item \textbf{CLIQUE} $\leq^p$ \textbf{ISO\_SUBGRAPH}, es decir, se puede reducir el problema de \textbf{CLIQUE} a \textbf{ISO\_SUBGRAPH} en tiempo polinómico.
    \end{itemize}
\end{frame}

\section{ISO\_SUBGRAPH está en NP}
\begin{frame}{ISO\_SUBGRAPH está en NP}
    Para demostrar que \textbf{ISO\_SUBGRAPH} está en NP, debemos comprobar que podemos verifcar una solución en tiempo polinómico.

    Debemos verificar que existe un subgrafo $G'$ de $G$ y una biyección $f : V_{G'} \to V_H$ que preserve las aristas. Esto implica verificar lo siguiente:

    \begin{itemize}
        \item $G'$ es un subgrafo de $G$ (comprobar que $V_{G'} \subseteq V_G$ y que $A_{G'} = A_G \cap (V_{G'} \times V_{G'})$).
        \item $f$ es una biyección entre $V_{G'}$ y $V_H$.
        \item Para cada arista $(u, v) \in A_{G'}$, verificar que $(f(u), f(v)) \in A_H$ y viceversa.
    \end{itemize}

    Dado que todas estas verificaciones se pueden realizar en tiempo
    polinómico, y en cantidades polinómicas, concluimos que \textbf{ISO\_SUBGRAPH} está en NP.
\end{frame}

\section{Reducción de CLIQUE a ISO\_SUBGRAPH}
\begin{frame}{Reducción de CLIQUE a ISO\_SUBGRAPH}
    Para demostrar que \textbf{CLIQUE} $\leq^p$ \textbf{ISO\_SUBGRAPH}, debemos construir una reducción polinómica de cualquier instancia de \textbf{CLIQUE} a una instancia de \textbf{ISO\_SUBGRAPH}.

    Dado un grafo $G$ y un entero $k$, queremos determinar si $G$ tiene una clique de tamaño $k$. Para esto, construiremos un grafo $H$ que sea un clique de tamaño $k$, es decir, $H = (V_H, A_H)$ donde $V_H = \{1, 2, \ldots, k\}$ y $A_H = \{(i, j) : 1 \leq i < j \leq k\}$.

    Luego, la instancia de \textbf{ISO\_SUBGRAPH} será el par $(G, H)$. La pregunta es si $G$ tiene un subgrafo isomorfo a $H$, lo cual es equivalente a preguntar si $G$ tiene una clique de tamaño $k$.

    Por lo tanto, si podemos resolver \textbf{ISO\_SUBGRAPH} para esta instancia, podremos resolver \textbf{CLIQUE}. Esto muestra que \textbf{CLIQUE} se reduce a \textbf{ISO\_SUBGRAPH} en tiempo polinómico.
\end{frame}

% Sección 3
\section{Conclusión}
\begin{frame}{Conclusión}
    Hemos demostrado que \textbf{ISO\_SUBGRAPH} es NP-completo al mostrar que:

    \begin{itemize}
        \item \textbf{ISO\_SUBGRAPH} está en NP, ya que podemos verificar una solución en tiempo polinómico.
        \item \textbf{CLIQUE} se reduce a \textbf{ISO\_SUBGRAPH} en tiempo polinómico, lo que implica que si pudiéramos resolver \textbf{ISO\_SUBGRAPH} eficientemente, también podríamos resolver \textbf{CLIQUE}.
    \end{itemize}

    Por lo tanto, concluimos que \textbf{ISO\_SUBGRAPH} es un problema NP-completo.
\end{frame}

\end{document}